% SOURCE_AUTHORITY: sga5_fr_workpass.tex, lines 11309--11348 (LF-normalized unit).
% SOURCE_SHA256: 791F4EFFC5E02832D5D77ED03518C8156D6F07E4C8238B03545DB93D883FBB28
\emph{Demonstração.} Demonstremos (i). Ponhamos $\xi=c_1(\mathcal O_{Y'}(1))=k^*(\bar\xi)$. Repetindo o argumento da demonstração de (8.4.1), obtêm-se as igualdades
\[
c_d(\check E)=\sum_{p=0}^d g_1^*(c_p(\check N))\bar\xi^{\,d-p},
\tag{9.7.1}
\]
\[
c_{d-1}(\check F)=\sum_{q=0}^{d-1}g^*(c_q(\check N))\xi^{d-1-q}.
\tag{9.7.2}
\]
De (9.7.2) deduz-se
\[
k_*[c_{d-1}(\check F)]
=\sum_{q=0}^{d-1}k_*k^*[g_1^*(c_q(\check N))\bar\xi^{\,d-1-q}],
\]
de onde, graças à fórmula de projeção para $k$ e ao fato de que $k_*(1)=\bar\xi$,
\[
k_*[c_{d-1}(\check F)]=\sum_{q=0}^{d-1}g_1^*(c_q(\check N))\bar\xi^{\,d-q}.
\tag{9.7.3}
\]
A afirmação (i) resulta imediatamente da comparação (9.7.1) e (9.7.3). Demostremos (ii). Pode-se supor que $y$ é um ciclo irredutível de $Y$, correspondente a uma subvariedade denotada da mesma forma. Então o ciclo $\gamma_*(y)=y\times0$ é racionalmente equivalente a $y\times1$, então basta ver a igualdade de ciclos
\[
\beta_*f_2^*(y\times1)=f_1^*\alpha_*(y\times1).
\tag{9.7.4}
\]
Mas, como $(y\times1)\cap(Y\times0)=\emptyset$, $f_1$ é um isomorfismo sobre um ambiente aberto de $y\times1$, de onde resulta imediatamente (9.7.4). Demostremos (iii). Denotando por $p:X\times\mathbb P^1\to X$ a primeira projeção, resulta de (9.1.2) que
\[
u_*(t)=p^*(t)\cdot u_*(1),
\tag{9.7.5}
\]
de onde, comparando com (9.3.4),
\[
f_1^*u_*(t)=(p\circ f_1)^*(t)\bigl[(j_1)_*(1)+v_*(1)\bigr].
\tag{9.7.6}
\]
Utilizando as fórmulas de projeção para $j_1$ e $v$, respectivamente, deduz-se
\[
f_1^*u_*(t)=(j_1)_*\bigl[(p\circ f_1\circ j_1)^*(t)\bigr]+v_*\bigl[(p\circ f_1\circ v)^*(t)\bigr],
\tag{9.7.7}
\]
de onde (iii), já que $p\circ f_1\circ j_1=(p\circ u)\circ i\circ g_1=i\circ g_1$ e $p\circ f_1\circ v=(p\circ u)\circ f=f$.
